\section{Лекция 2. Отношения и отображения}

\subsection{Бинарные отношения}

\begin{definition}
    Говорят, что между элементами множеств \(X\) и \(Y\) задано \textdef{бинарное отношение} \(R\), если указано подмножество их декартова произведения:
    \[
        R \subseteq X \times Y
    \]
\end{definition}

\noindent Частный случай: когда \(X = Y\), отношение задаётся как \(R \subseteq X^2\).

\begin{designation}
    Тот факт, что пара \((x, y)\) принадлежит отношению \(R\) (то есть \((x, y) \in R\)), часто записывают инфиксным образом: \(x R y\). В зависимости от контекста используются символы: \(x \preceq y\), \(x \sim y\), \(x \prec y\) и т.д.
\end{designation}

\begin{definition}
    Для отношения \(R \subseteq X \times Y\) определяются:
    \begin{itemize}
        \item \textdef{Область определения} (domain): \(\operatorname{dom} R = \{x \in X \mid \exists y \in Y: x R y\}\);
        \item \textdef{Область значений} (image): \(\operatorname{im} R = \{y \in Y \mid \exists x \in X: x R y\}\).
    \end{itemize}
\end{definition}

\subsubsection{Свойства отношений на множестве}
Рассмотрим отношение \(R\) на одном множестве (\(R \subseteq X^2\)). Основные свойства, которыми оно может обладать:

\begin{enumerate}
    \item \textdef{Рефлексивность:} \(\forall x \in X \implies x R x\);
    \item \textdef{Симметричность:} \(\forall x, y \in X : x R y \implies y R x\);
    \item \textdef{Антисимметричность:} \(\forall x, y \in X : (x R y \ \text{и} \ y R x) \implies x = y\) (характерно для нестрогих порядков, \(\preceq\));
    \item \textdef{Асимметричность:} \(\forall x, y \in X : x R y \implies \neg(y R x)\) (характерно для строгих порядков, \(\prec\));
    \item \textdef{Транзитивность:} \(\forall x, y, z \in X : (x R y \ \text{и} \ y R z) \implies x R z\).
\end{enumerate}

\subsection{Отношения эквивалентности}

\begin{definition}
    Отношение \(R\) называется \textdef{отношением эквивалентности} (обозначается \(\sim\)), если оно одновременно рефлексивно, симметрично и транзитивно.
\end{definition}

\begin{example}
    Приведём примеры:
    \begin{enumerate}
        \item \emph{Полное отношение:} \(x R y\) для всех \(x, y \in X\) (все элементы эквивалентны друг другу).
        \item \emph{Отношение равенства (тождества):} никакие различные \(x\) и \(y\) не связаны, каждый элемент связан только сам с собой.
    \end{enumerate}
\end{example}

\begin{definition}
    \textdef{Класс эквивалентности} элемента \(a \in X\) — это множество всех элементов из \(X\), эквивалентных \(a\):
    \[
        \operatorname{cl}_\sim(a) = \{x \in X \mid x \sim a\}
    \]
\end{definition}

\begin{theorem}
    Класс эквивалентности определяется любым своим представителем:
    \[
        b \in \operatorname{cl}_\sim(a) \implies \operatorname{cl}_\sim(b) = \operatorname{cl}_\sim(a)
    \]
\end{theorem}

\begin{proof}
    Пусть \(b \in \operatorname{cl}_\sim(a)\), то есть \(b \sim a\). Покажем взаимное включение множеств.
    
    1. Пусть \(x \in \operatorname{cl}_\sim(b)\), тогда \(x \sim b\). Так как \(b \sim a\), по свойству транзитивности получаем \(x \sim a\), значит \(x \in \operatorname{cl}_\sim(a)\). Следовательно, \(\operatorname{cl}_\sim(b) \subseteq \operatorname{cl}_\sim(a)\).
    
    2. Пусть \(x \in \operatorname{cl}_\sim(a)\), тогда \(x \sim a\). Так как \(b \sim a\), по свойству симметричности \(a \sim b\). Тогда из \(x \sim a\) и \(a \sim b\) по транзитивности следует \(x \sim b\), значит \(x \in \operatorname{cl}_\sim(b)\). Следовательно, \(\operatorname{cl}_\sim(a) \subseteq \operatorname{cl}_\sim(b)\).
    
    Из двух включений следует равенство: \(\operatorname{cl}_\sim(a) = \operatorname{cl}_\sim(b)\).
\end{proof}

\begin{theorem}
    Любые два класса эквивалентности либо не пересекаются, либо полностью совпадают.
\end{theorem}

\begin{proof}
    Пусть два класса \(\operatorname{cl}_\sim(a)\) и \(\operatorname{cl}_\sim(b)\) имеют хотя бы один общий элемент: \(x \in \operatorname{cl}_\sim(a) \cap \operatorname{cl}_\sim(b)\). 
    Это означает, что \(x \sim a\) и \(x \sim b\). Из симметричности \(a \sim x\). Из транзитивности (\(a \sim x\) и \(x \sim b\)) получаем \(a \sim b\).
    А значит, \(a \in \operatorname{cl}_\sim(b)\). По предыдущей теореме это означает, что \(\operatorname{cl}_\sim(a) = \operatorname{cl}_\sim(b)\).
\end{proof}

\begin{corollary}
    Отношение эквивалентности разбивает исходное множество на непересекающиеся классы эквивалентности (формирует разбиение множества).
\end{corollary}

\begin{definition}
    Множество всех классов эквивалентности множества \(X\) по отношению \(\sim\) называется \textdef{фактор-множеством} и обозначается \(X/{\sim}\) (в конспектах также встречается \(|X|_\sim\)).
\end{definition}

\begin{remark}
    Переход от элементов к их классам называется \textdef{факторизацией}. Это мощный инструмент, позволяющий, например, сводить работу с бесконечными множествами к работе с конечными фактор-множествами.
\end{remark}

\subsection{Отношения порядка}

\begin{definition}
    \textdef{Отношением частичного порядка} (нестрогого) называется отношение, которое рефлексивно, антисимметрично и транзитивно.
\end{definition}

\begin{example}
    \(\le\) для чисел, \(\subseteq\) для множеств.
\end{example}

\begin{definition}
    \textdef{Отношением строгого частичного порядка} называется отношение, которое иррефлексивно, асимметрично и транзитивно.
\end{definition}

\begin{example}
     \(<\) для чисел, \(\subset\) для множеств.
\end{example}

\vspace{1 em}

\begin{remark}
    Слово «частичный» означает, что не любые два элемента обязаны быть сравнимы. Существуют элементы, для которых не выполняется ни \(x \preceq y\), ни \(y \preceq x\).
\end{remark}

\vspace{1 em}

\begin{remark}
    
    В математике само по себе выражение \((x_1,y_1) \preceq (x_2,y_2)\) не имеет
    однозначного смысла, пока не указано, какой порядок на парах/кортежах
    имеется в виду. Чаще всего используются два порядка:

    \begin{quote}
        
        1. \textbf{Покоординатный порядок.} Считаем, что:
        \[
            (x_1,y_1) \preceq (x_2,y_2)
            \quad\Longleftrightarrow\quad
            x_1 \le x_2 \ \text{ и } \ y_1 \le y_2.
        \]
        То есть сравнение идёт по каждой координате отдельно.
        
        Это частичный порядок: не любые две пары сравнимы.
        Например, \((1,5)\quad\text{и}\quad(2,3)\) не сравнимы, потому
        что \(1 < 2\), но \(5 > 3\).

        Если пишут строго:
        \[
            (x_1,y_1) \prec (x_2,y_2),
        \]
        то обычно имеют в виду
        \[
            (x_1,y_1) \preceq (x_2,y_2)
            \quad\text{и}\quad
            (x_1,y_1) \ne (x_2,y_2).
        \]

        2. \textbf{Лексикографический порядок.} Так сравнивают слова в словаре. Сначала
        сравниваем первые координаты; если они равны, то сравнивают вторые:
        \[
            (x_1,y_1) \preceq_{\mathrm{lex}} (x_2,y_2)
            \quad\Longleftrightarrow\quad
            x_1 < x_2
            \ \text{ или }
            \bigl(x_1 = x_2 \ \text{ и } \ y_1 \le y_2\bigr).
        \]
        Этот порядок линейный: любые две пары сравнимы. Пример:
        \((1,10) \prec_{\mathrm{lex}} (2,0),\) потому что \(1 < 2\), хотя
        вторая координата больше.

    \end{quote}

\end{remark}

\begin{definition}
    Пусть задано частично упорядоченное множество (ЧУМ) \(\langle X, \preceq \rangle\).
    \begin{itemize}
        \item Элемент \(x_{\min} \in X\) называется \textdef{минимальным}, если не существует элемента строго меньшего него. То есть \(\forall x \in X\): либо \(x\) и \(x_{\min}\) несравнимы, либо \(x_{\min} \preceq x\).
        \item Элемент \(x_{\text{наим}} \in X\) называется \textdef{наименьшим}, если он меньше или равен любому другому элементу: \(\forall x \in X \implies x_{\text{наим}} \preceq x\).
    \end{itemize}
    Минимальных элементов в множестве может быть несколько (или не быть вовсе). Наименьший элемент (если он существует) всегда единственен.
\end{definition}

\begin{definition}
    Частично упорядоченное множество \(\langle X, \preceq \rangle\) называется \textdef{фундированным}, если в любом его непустом подмножестве существует минимальный элемент.
    \\ Эквивалентное определение: в множестве не существует бесконечно убывающих цепей \(x_1 \succ x_2 \succ x_3 \succ \dots\)
\end{definition}

\begin{definition}
    Множество называется \textdef{вполне упорядоченным}, если на нём задан линейный порядок (любые два элемента сравнимы) и оно является фундированным.
\end{definition}

\begin{example}
    Приведём примеры:

    \(\langle \N, \le \rangle\) --- фундировано (вполне упорядочено).
    В нём любое непустое подмножество имеет наименьший элемент,
    и нельзя бесконечно убывать:
    \(5 > 4 > 3 > 2 > 1 > 0\) --- дальше некуда.

    \(\langle \Z, \le \rangle\) --- не фундировано (есть бесконечно
    убывающая цепь). Например, подмножество всех отрицательных чисел
    не имеет минимального элемента, и можно бесконечно убывать:
    \(0 > -1 > -2 > -3 > \ldots\)
\end{example}


\subsection{Отображения}

\begin{definition}
    Бинарное отношение \(f \subseteq X \times Y\) называется \textdef{отображением} (или функцией), если каждому элементу из \(X\) ставится в соответствие не более одного элемента из \(Y\):
    \[
        (x f y \ \text{и} \ x f z) \implies y = z \quad \forall x \in X, \ \forall y, z \in Y
    \]
    Обозначение: \(f: X \to Y\), \(y = f(x)\).
\end{definition}

\begin{definition}
    Отображение \(f: X \to Y\) называется \textdef{полным}, если \(\operatorname{dom} f = X\), и \textdef{частичным} в противном случае. \emph{Далее будут подразумеваться только полные отображения.}
\end{definition}

\subsubsection{Свойства отображений}
\begin{enumerate}
    \item \textdef{Инъективное:} разные элементы переходят в разные. \(f(x_1) \neq f(x_2)\), если \(x_1 \neq x_2\) (или эквивалентно: \(f(x_1) = f(x_2) \implies x_1 = x_2\)).
    \item \textdef{Сюръективное:} каждый элемент из \(Y\) имеет прообраз. \(\forall y \in Y \ \exists x \in X : f(x) = y\).
    \item \textdef{Биективное:} одновременно инъективно и сюръективно (взаимно однозначное).
\end{enumerate}

\subsubsection{Композиция и обратные отображения}

\begin{definition}
    \textdef{Произведением (композицией)} отображений \(g: X \to Y\) и \(f: Y \to Z\) называется отображение \(f \circ g: X \to Z\), действующее по правилу: \((f \circ g)(x) = f(g(x))\).
\end{definition}

\begin{theorem}
    Произведение (композиция) отображений ассоциативна. То есть для любых
    отображений \(h\), \(f\) и \(g\) верно следующее:
    \[
        h \circ ( f \circ g) = (h \circ f) \circ g
    \]
\end{theorem}

\begin{proof}
    Возьмём произвольный элемент \(x\) из области определения \(g\).

    \begin{quote}
        Тогда
        \[
            \bigl(h \circ (f \circ g)\bigr)(x)
            = h\bigl((f \circ g)(x)\bigr)
            = h\bigl(f(g(x))\bigr),
        \]
        и точно так же
        \[
            \bigl((h \circ f) \circ g\bigr)(x)
            = (h \circ f)\bigl(g(x)\bigr)
            = h\bigl(f(g(x))\bigr).
        \]
        Обе композиции дают одно и то же значение \(h(f(g(x)))\) в каждой точке \(x\).
        Значит, отображения равны.
    \end{quote}
\end{proof}

\begin{definition}
    Для отображения \(f: X \to Y\) отображение \(f^{-1}: Y \to X\) называется \textdef{обратным}, если выполняются условия:
    \[
        \begin{cases}
            f \circ f^{-1} = e_Y \\
            f^{-1} \circ f = e_X
        \end{cases}
    \]
    где \(e_X, e_Y\) --- тождественные отображения на множествах \(X\) и \(Y\) соответственно (\(e(x) = x\)).
\end{definition}

\begin{remark}
    Далеко не каждое отображение имеет обратное.
\end{remark}

\begin{lemma} \label{lem:inv_inj_surj}
    Пусть даны отображения \(f: X \to Y\) и \(g: Y \to X\), причём \(g \circ f = e_X\). Тогда \(f\) — инъективно, а \(g\) — сюръективно.
\end{lemma}

\begin{proof}
    Докажем поэтапно.
    \begin{quote}

        1. \textbf{Инъективность \(f\)} докажем от противного. Пусть \(f\) не инъективно, то есть \(\exists x_1, x_2 \ (x_1 \neq x_2)\), такие что \(f(x_1) = f(x_2)\). Применим к обеим частям отображение \(g\):
        \[
            g(f(x_1)) = g(f(x_2)) \implies (g \circ f)(x_1) = (g \circ f)(x_2)
        \]
        Так как по условию \(g \circ f = e_X\), получаем \(e_X(x_1) = e_X(x_2)\), то есть \(x_1 = x_2\). Противоречие. Значит, \(f\) инъективно.

        2. \textbf{Сюръективность \(g\)} докажем напрямую. Нам нужно показать, что \(\forall x \in X \ \exists y \in Y : g(y) = x\).
        Возьмём произвольный \(x \in X\) и положим \(y = f(x)\). Тогда \(g(y) = g(f(x)) = (g \circ f)(x) = e_X(x) = x\). Такой \(y\) найден, следовательно, \(g\) сюръективно.
    
    \end{quote}

\end{proof}

\begin{theorem}[Критерий обратимости]
    Отображение \(f: X \to Y\) обратимо тогда и только тогда, когда оно биективно.
\end{theorem}

\begin{proof}
    Докажем в обе стороны.

    \begin{quote}
        \(\implies\) \textbf{Обратимость влечет биективность:}
        Если \(f\) обратимо, существует \(f^{-1}\). Из равенства \(f^{-1} \circ f = e_X\) по Лемме \ref{lem:inv_inj_surj} следует инъективность \(f\). Из равенства \(f \circ f^{-1} = e_Y\) по той же Лемме следует сюръективность \(f\). Таким образом, \(f\) биективно.
        
        \(\impliedby\) \textbf{Биективность влечет обратимость:}
        Так как \(f\) сюръективно, \(\forall y \in Y \ \exists x \in X : f(x) = y\). Так как \(f\) инъективно, такой \(x\) единственен. Определим отображение \(f^{-1}(y) = x\). Проверим композиции:
        \[
            (f \circ f^{-1})(y) = f(f^{-1}(y)) = f(x) = y = e_Y(y) \implies f \circ f^{-1} = e_Y
        \]
        \[
            (f^{-1} \circ f)(x) = f^{-1}(f(x)) = f^{-1}(y) = x = e_X(x) \implies f^{-1} \circ f = e_X
        \]
        Отображение \(f^{-1}\) существует и удовлетворяет определению обратного.
    \end{quote}

\end{proof}


\begin{theorem}
    \emph{Свойства композиции:}

    \begin{enumerate}
        \item Композиция инъективных отображений --- инъективна.
        \item Композиция сюръективных отображений --- сюръективна.
        \item Композиция биективных отображений --- биективна.
    \end{enumerate}

\end{theorem}

\begin{proof}
    Пусть \(g: X \to Y\) и \(f: Y \to Z\).
    \begin{quote}
        Рассмотрим \(f \circ g: X \to Z\).
        
        1. От противного. Пусть \(\exists x_1 \neq x_2\), такие что \(f(g(x_1)) = f(g(x_2))\). Так как \(f\) инъективно, аргументы равны: \(g(x_1) = g(x_2)\). Так как \(g\) инъективно, \(x_1 = x_2\). Противоречие.
        
        2. Нужно показать, что \(\forall z \in Z \ \exists x \in X : f(g(x)) = z\). Так как \(f\) сюръективно, для \(z\) найдется \(y \in Y\), что \(f(y) = z\). Так как \(g\) сюръективно, для этого \(y\) найдется \(x \in X\), что \(g(x) = y\). Отсюда \(f(g(x)) = f(y) = z\).
        
        3. Следствие из пунктов 1 и 2.
    \end{quote}
\end{proof}

\begin{theorem}
    Если \(f\) и \(g\) обратимы, то их композиция \(f \circ g\) также обратима, причём:
    \[
        (f \circ g)^{-1} = g^{-1} \circ f^{-1}
    \]
\end{theorem}

\begin{proof}
    Обозначим \(h = g^{-1} \circ f^{-1}\).

    \begin{quote}
        Проверим свойства обратного отображения для \(f \circ g\), используя ассоциативность композиции:
        \begin{align*}
            (f \circ g) \circ h &= (f \circ g) \circ (g^{-1} \circ f^{-1}) = f \circ (g \circ g^{-1}) \circ f^{-1} = f \circ e_Y \circ f^{-1} = f \circ f^{-1} = e_Z \\
            h \circ (f \circ g) &= (g^{-1} \circ f^{-1}) \circ (f \circ g) = g^{-1} \circ (f^{-1} \circ f) \circ g = g^{-1} \circ e_Y \circ g = g^{-1} \circ g = e_X
        \end{align*}
        Следовательно, \(h\) действительно является обратным к \(f \circ g\).
    \end{quote}
\end{proof}


\begin{theorem}
    Обратное отображение определяется однозначно.
\end{theorem}

\begin{proof}
    Докажем, предположив обратное.

    \begin{quote}
        Пусть у отображения \(f: X \to Y\) есть два обратных отображения: \(g_1\) и \(g_2\). По определению обратного, \(f \circ g_2 = e_Y\).
        Запишем цепочку равенств:
        \[
            g_1 = g_1 \circ e_Y = g_1 \circ (f \circ g_2) = (g_1 \circ f) \circ g_2 = e_X \circ g_2 = g_2
        \]
        Следовательно, \(g_1 = g_2\).
    \end{quote}
\end{proof}

\subsubsection{Алгебраические структуры отображений}

Тождественное отображение играет роль нейтрального элемента: \(f \circ e_X = e_Y \circ f = f\).

Если рассматривать множество всех отображений из множества в само себя (\(f: X \to X\)), мы получаем следующие алгебраические структуры:
\begin{enumerate}
    \item \textdef{Полугруппа} — множество отображений с ассоциативной операцией композиции.
    \item \textdef{Моноид} — полугруппа, в которой есть нейтральный элемент (\(e_X\)).
    \item \textdef{Группа} — если рассматривать \emph{только биективные} отображения \(X \to X\), то каждое из них имеет обратное (\(f^{-1}\)). Моноид с обратимыми элементами образует группу.
\end{enumerate}